\chapter{全书总览与使用路径}

\section{本手册的工作方式}

这本手册按“题目识别”而不是按“算法教材目录”组织。考场上你通常不是先想起某个算法，再去找题目匹配它；更常见的是先看到题面中的若干信号，例如区间、多次询问、最短、连通、排序、字符串、日期、进制、命令序列，然后再判断应该翻到哪个章节。

因此，本手册的每个决策入口都尽量回答下面七个问题：

\begin{enumerate}
  \item 题面中出现了什么信号？
  \item 数据范围允许什么复杂度？
  \item 它更像 A、B、C、D、E 中哪类题？
  \item 先尝试哪个模板？
  \item 需要使用哪些 C++ 语法或 STL 容器？
  \item 最容易错在哪里？
  \item 如果满分做法不会，能否先写部分分？
\end{enumerate}

\Warn{使用原则：}不要在考场上临时学习一个完全没练过的模板。模板页是用来提醒和防错的，不是用来从零学习的。

\section{全书结构路线图}

本书按考场使用顺序重排为四个分部。先用索引定位，再用题型章节展开，最后用模板和真题反向校准。

\begin{longtable}{p{3.2cm}p{4.2cm}p{6.6cm}}
\toprule
分部 & 章节 & 用法 \\
\midrule
导航与判断 & 第 1--3 章 & 先看全书用法，再用 1 分钟极速索引定位，必要时进入详细判断树。 \\
基础与题型 & 第 4--8 章 & 先补 C++ 和 STL 基础，再按 A/B/C/D/E 题型查看策略和高频入口。 \\
模板与检查 & 第 9--10 章 & 找可复制模板，并在提交前扫易错点。 \\
真题校准 & 第 11 章 & 用历年题反查模板和判断树，发现误判后回到第 3 章修正。 \\
\bottomrule
\end{longtable}

推荐阅读路线：

\begin{enumerate}
  \item \Key{考前通读：}第 1 章、第 4 章、第 5--8 章、第 10 章。
  \item \Key{刷题复盘：}第 3 章判断路径、第 9 章模板、第 11 章真题索引。
  \item \Key{考场翻查：}第 2 章极速索引；定位不准时翻第 3 章；写代码时翻第 9 章；提交前翻第 10 章。
\end{enumerate}

\section{官方规则速查}

本节只作为考前检查清单。正式打印前必须重新访问 CCF CSP 官网核对最新版本。

\begin{longtable}{p{3.2cm}p{10.8cm}}
\toprule
项目 & 当前应记录的信息 \\
\midrule
官方网站 & \url{https://www.cspro.org/} \\
认证形式 & 上机编程，黑盒测试，按测试点得分。 \\
题目数量 & 通常为 5 道编程题。每题 100 分，总分 500 分。 \\
考试时间 & 通常为 240 分钟。 \\
输入输出 & 使用标准输入、标准输出。不要输出提示语、调试信息或多余空格。 \\
提交方式 & 每题提交一个源文件；同一题多次提交时，一般以最后一次提交为准。 \\
语言环境 & 以官网报名系统和当次考试公告为准。C++ 版本、编译器和操作系统可能变化。 \\
纸质资料 & 根据 CCF 官网答疑和考生须知，应考时可携带纸质书籍、打印资料、手写资料；不得使用通讯工具、存储设备。考前必须再次核对。 \\
程序长度 & 官网注意事项中提到程序长度限制，准备模板时不要复制大量无关代码。 \\
\bottomrule
\end{longtable}

建议在打印版首页手写确认：

\begin{itemize}
  \item 我已经在考前一周核对了官网规则。
  \item 我已经确认了本次考试允许携带的纸质资料范围。
  \item 我已经确认了本次考试 C++ 编译标准。
  \item 我知道不能携带电子资料、存储设备或联网设备。
\end{itemize}

\section{考场总流程}

\subsection{开场 10 分钟}

开场不要立刻写代码。先快速浏览 5 道题，目标是建立时间分配和得分路线。

\begin{enumerate}
  \item 看每道题的题面长度。题面短不一定简单，题面长往往意味着模拟或细节多。
  \item 看每道题的数据范围。数据范围比题面故事更重要。
  \item 判断 A、B 是否可以快速完成。
  \item 判断 C 是算法题还是大模拟。
  \item 判断 D、E 是否有子任务、小数据、特殊性质或熟悉模板。
  \item 在草稿纸上写下计划：先做哪题、卡住多久换题、最后回来看哪题。
\end{enumerate}

\subsection{每题固定流程}

每道题都按同一套流程处理：

\begin{enumerate}
  \item 读题面，圈出关键词。
  \item 抄下所有数据范围。
  \item 估算可接受复杂度。
  \item 查本章总决策树，定位章节。
  \item 写出暴力思路，用来理解题意和造样例。
  \item 如果满分模板明确，再写正式代码。
  \item 先过样例，再造边界样例。
  \item 删除调试输出后提交。
\end{enumerate}

\Warn{不要跳过复杂度估算。}很多新手失分不是因为不会算法，而是把 $O(n^2)$ 写到了 $n=10^5$ 的题里。

\subsection{建议时间分配}

\begin{longtable}{p{2cm}p{3cm}p{5cm}p{4cm}}
\toprule
题号 & 建议时间 & 目标 & 卡住时的处理 \\
\midrule
A & 20--35 分钟 & 稳定满分 & 15 分钟没思路就重新读题；30 分钟没过样例先跳。 \\
B & 40--70 分钟 & 尽量满分 & 优先查前缀和、差分、排序、哈希、BFS。 \\
C & 60--100 分钟 & 高分或满分 & 如果是大模拟，先拆结构体和函数，不要全写在主函数。 \\
D & 40--80 分钟 & 部分分到满分 & 先找小数据和熟悉模板。不会满分也要写暴力。 \\
E & 剩余时间 & 尽量拿分 & 重点看子任务、特殊性质、样例规律。 \\
\bottomrule
\end{longtable}

时间分配不是固定规则。若 A、B 很顺，可以把更多时间留给 C、D；若 A、B 卡住，不要让前两题吞掉整场考试。

\section{数据范围与复杂度}

\subsection{复杂度速查表}

\begin{longtable}{p{4cm}p{5cm}p{5cm}}
\toprule
输入规模 & 通常可考虑 & 常见模板 \\
\midrule
$n \le 20$ & 指数搜索、状态压缩、全排列、DFS 暴力 & DFS、状压 DP、枚举子集 \\
$n \le 100$ & $O(n^3)$ 通常可试 & Floyd、区间 DP、三重枚举 \\
$n \le 500$ & $O(n^3)$ 可能可行，需看时限 & Floyd、部分 DP \\
$n \le 1000$ & $O(n^2)$ 通常可行 & 双重循环、二维 DP、朴素图算法 \\
$n \le 5000$ & $O(n^2)$ 需谨慎 & 优化前的 DP、枚举配合剪枝 \\
$n \le 10^5$ & $O(n \log n)$ 或 $O(n)$ & 排序、二分、前缀和、差分、堆、树状数组 \\
$n \le 10^6$ & $O(n)$ 或较轻的 $O(n \log n)$ & 线性扫描、计数、前缀和 \\
$n \ge 10^7$ & 通常只能线性或公式 & 数学、简单扫描、避免大对象数组 \\
\bottomrule
\end{longtable}

\subsection{复杂度校验规则}

\begin{itemize}
  \item 看到双重循环，先计算最坏次数。
  \item 有多组数据时，总复杂度要乘以组数，或者看所有 $n$ 的总和限制。
  \item 排序一般是 $O(n \log n)$，$n=10^5$ 常见可行。
  \item Floyd 是 $O(n^3)$，只适合点数较小的图。
  \item Dijkstra 配合堆是 $O((n+m)\log n)$，适合正权稀疏图。
  \item BFS 适合边权全为 1 或每步代价相同的最短步数。
  \item 前缀和适合不修改的区间查询。
  \item 差分适合批量区间修改后统一查询。
  \item 树状数组和线段树适合修改与查询交织出现。
\end{itemize}

\Warn{溢出提醒：}只要出现求和、乘法、路径长度、方案数、时间累计，默认先考虑 \Code{long long}。两个 \Code{int} 相乘时写成 \Code{1LL * a * b}。

\section{A/B/C/D/E 总体策略}

\subsection{A 题}

A 题通常考察基本实现能力。它可能是简单模拟、统计、字符串、日期、进制、排序或公式题。

\begin{itemize}
  \item 目标：尽量一次 AC。
  \item 优先检查：输入格式、输出格式、整数溢出、小数精度、边界值。
  \item 常用章节：C++ 基础、数组统计、字符串、进制、日期、排序。
  \item 不建议：为了追求代码短而写不熟悉的技巧。
\end{itemize}

\subsection{B 题}

B 题常常是第一个真正需要算法识别的题，但算法通常不会太深。

\begin{itemize}
  \item 高频方向：排序、前缀和、差分、哈希、二分、BFS、简单图、简单 DP。
  \item 读题重点：是否有多次询问，是否有区间，是否需要快速统计。
  \item 常见错误：用暴力过样例但过不了大数据。
\end{itemize}

\subsection{C 题}

C 题经常是分水岭。它可能是大模拟，也可能是中等算法题。

\begin{itemize}
  \item 如果题面很长，先判断是否是规则模拟。
  \item 如果输入是命令、事件、状态变化，优先设计结构体和函数。
  \item 如果出现图、网格、路径、依赖关系，再进入对应算法分支。
  \item 大模拟题不要把所有逻辑塞进 \Code{main}。
\end{itemize}

\subsection{D/E 题}

D、E 题要同时考虑满分算法和部分分。

\begin{itemize}
  \item 先看是否给出子任务或特殊数据范围。
  \item 先写能过小数据的暴力，避免空题。
  \item 若满分算法明显且练过，再冲满分。
  \item 若 20 分钟没有形成可写方案，转为部分分策略。
\end{itemize}

\section{总决策树}

\subsection{第一层：看题面信号}

\begin{longtable}{p{4.2cm}p{4.8cm}p{5cm}}
\toprule
题面信号 & 优先怀疑 & 先翻章节 \\
\midrule
字符、单词、路径、命令、表达式、括号 & 字符串处理或解析模拟 & 字符串专题、栈、大模拟 \\
二进制、十六进制、位、补码、编码 & 进制和位运算 & 进制专题、位运算 \\
日期、时间、星期、天数、日历 & 日期时间处理 & 日期专题 \\
区间和、多次查询、子矩阵和 & 前缀和 & 前缀和模板 \\
区间加、范围修改、最后输出 & 差分 & 差分模板 \\
修改和查询交替出现 & 树状数组或线段树 & 数据结构模板 \\
最大值、最小值动态维护 & 堆、集合、单调结构 & 优先队列、set、单调队列 \\
排序、排名、优先级、相同条件按某字段 & 自定义排序 & sort 与 lambda \\
去重、出现次数、频率统计 & 哈希或 map & map、unordered map、计数数组 \\
最少步数、迷宫、网格移动 & BFS & BFS、网格 BFS \\
正权最短路、道路费用、距离最小 & Dijkstra & 最短路模板 \\
任意两点最短路且点数小 & Floyd & Floyd 模板 \\
合并集合、是否连通、同一类 & 并查集 & 并查集模板 \\
依赖、先后顺序、任务调度 & 拓扑排序 & 拓扑排序模板 \\
最大化最小值、最小化最大值、答案可行性 & 二分答案 & 二分答案模板 \\
容量、价值、选择若干个、最多或恰好 & 背包 DP & 动态规划模板 \\
规则很多、对象很多、按时间变化 & 大模拟 & 大模拟专题 \\
\bottomrule
\end{longtable}

\subsection{第二层：用数据范围排除错误方向}

\begin{longtable}{p{5cm}p{4.5cm}p{4.5cm}}
\toprule
如果你想写 & 但数据范围是 & 应改查 \\
\midrule
枚举所有点对 $O(n^2)$ & $n=10^5$ & 排序、哈希、前缀和、树状数组、双指针 \\
枚举所有区间 $O(n^2)$ & $n=10^5$ & 前缀和、双指针、单调队列、DP 优化 \\
Floyd $O(n^3)$ & $n>500$ 或边很多 & Dijkstra、BFS、图的特殊性质 \\
每次查询都扫描数组 & 查询数 $q=10^5$ & 前缀和、树状数组、线段树、离线处理 \\
递归搜索所有方案 & $n>25$ 且无剪枝 & DP、贪心、数学、二分答案 \\
字符串反复插入删除 & 字符串长度和操作数很大 & 分块、链表思想、栈、模拟优化 \\
\bottomrule
\end{longtable}

\subsection{第三层：选择满分或部分分}

\begin{enumerate}
  \item 如果模板明确、数据范围匹配、自己练过，写满分做法。
  \item 如果模板大致明确但细节多，先写能跑的小版本，再逐步补全。
  \item 如果满分算法不明确，找子任务和小数据，写暴力或特殊情况。
  \item 如果题目很难但 A/B/C 还没稳，不要长时间停留在 D/E。
\end{enumerate}

\section{关键词到章节索引}

这张表用于快速翻书。后续章节完成后，需要把“章节编号”和“模板编号”补齐。

\begin{longtable}{p{4cm}p{4.5cm}p{5.2cm}}
\toprule
关键词 & 翻到哪里 & 核心提醒 \\
\midrule
字符串、字符、大小写 & 字符串专题 & 注意 \Code{cin} 和 \Code{getline} 混用。 \\
空格分割、命令行 & 字符串分割 & 用 \Code{stringstream} 或手写分割。 \\
括号、嵌套、表达式 & 栈、表达式解析 & 注意优先级和匹配失败。 \\
进制转换 & 进制专题 & 判断是否超过 \Code{long long}。 \\
二进制位 & 位运算 & 大位数用 \Code{1LL << k}。 \\
日期、星期 & 日期专题 & 先确认闰年和月份天数。 \\
区间查询 & 前缀和 & 不修改时优先前缀和。 \\
区间修改 & 差分 & 修改后统一查询时优先差分。 \\
在线修改查询 & 树状数组、线段树 & 看是单点还是区间。 \\
排序、排名 & sort 与 lambda & 相同条件要写完整。 \\
频率、出现次数 & map、unordered map、计数数组 & 值域小用数组更快。 \\
最短步数 & BFS & 边权相同才用 BFS。 \\
正权最短路 & Dijkstra & 距离用 \Code{long long}。 \\
连通、合并 & 并查集 & 路径压缩和按大小合并。 \\
依赖、先后 & 拓扑排序 & 入度为 0 入队。 \\
最优选择 & DP 或贪心 & 先看是否有状态转移。 \\
答案范围 & 二分答案 & 必须能写 \Code{check}。 \\
大量规则 & 大模拟 & 结构体、函数、状态表。 \\
\bottomrule
\end{longtable}

\section{低级错误总检查}

提交前至少检查下面内容：

\begin{itemize}
  \item 是否删除所有调试输出。
  \item 是否使用标准输入输出。
  \item 输出末尾是否多了不该有的内容。
  \item 是否把答案变量开成 \Code{long long}。
  \item 数组是否开到最大范围。
  \item 下标是从 0 还是从 1 开始。
  \item 循环边界是否包含末尾。
  \item 多组数据时是否清空数组、图、队列、map。
  \item 比较器是否满足严格弱序。
  \item \Code{lower\_bound} 或 \Code{find} 的返回值是否检查。
  \item \Code{priority\_queue} 是否需要小根堆。
  \item BFS 是否在入队时标记访问。
  \item Dijkstra 是否跳过过期状态。
  \item 递归深度是否可能爆栈。
\end{itemize}

\section{部分分总策略}

当满分做法暂时不会时，按这个顺序找分：

\begin{enumerate}
  \item \Key{小数据暴力：}如果 $n \le 1000$ 的测试点存在，写 $O(n^2)$；如果 $n \le 100$，考虑 $O(n^3)$。
  \item \Key{特殊结构：}图是树、链、完全图、DAG、无权图时，可能有更简单做法。
  \item \Key{特殊参数：}如果某个参数为 0、1、很小，单独写分支。
  \item \Key{无修改版本：}有修改和查询时，若部分数据没有修改，先用前缀和或静态算法。
  \item \Key{少量查询或少量修改：}查询少就每次暴力；修改少就重算或离线处理。
  \item \Key{样例和明显边界：}最后时间不足时，也要保证样例和简单边界正确。
\end{enumerate}

\Warn{部分分代码也要干净。}不要写一堆互相冲突的特判。特判越多，越容易把能拿的分弄丢。

\section{本章完成后的后续连接}

后续章节需要逐步补齐以下链接：

\begin{itemize}
  \item 将“关键词到章节索引”中的章节编号补全。
  \item 将每个决策树叶子节点连接到具体模板编号。
  \item 将真题索引中的每道题反向连接到本章的某个判断路径。
  \item 将个人错题中的错误补入“低级错误总检查”。
\end{itemize}
